%%%%

%%%   Самарский государственный технический университет

%%%   Кафедра прикладной математики и информатики

%%%

%%%   Преамбула для статьи в журнал "Вестник Самарского государственного

%%%   технического университета. Серия: «Физико-математические науки»"

%%%   Ориентирована под MikTEx 2.4 for Win32

%%%

%%%   Хотя сам журнал собирается под GNU/Linux на дистрибутиве TeXLive



\documentclass[11pt,a4paper,twoside]{article}



\usepackage[cp1251]{inputenc}

\usepackage[T2A]{fontenc}

\usepackage[english,russian]{babel}

\usepackage{samgtubib}

\usepackage[dvips]{graphicx}

\usepackage[multiple]{footmisc}



\usepackage{samgtu}

\renewcommand{\VestnikYear}{2009} % Год издания Вестника

\renewcommand{\VestnikNumber}{2\,(19)} % Номер Вестника

\renewcommand{\VestnikMonth}{Ноябрь} % Месяц выхода Вестника

\renewcommand{\VestnikData}{01.11.09} % Дата подписания Вестника в печать

\renewcommand{\VestnikUPL}{26,64} % Усл. печ. л.

\renewcommand{\VestnikUIL}{25,73} % Уч.-изд. л.

\renewcommand{\VestnikRegNumber}{479/09} % Регистрационный номер

\renewcommand{\VestnikOrder}{994} % Номер заказа






%
%
%
%\begin{document}
%
%
%\Partition{Физическая химия}{}

\renewcommand{\firstpage}{170}
\renewcommand{\lastpage}{179}

\udk{544.022:544.023.23}
\msc{92E10}

\title{Особенности структуры и кинетика роста пл\"eнок из адамантана в~тлеющем разряде}
{Особенности структуры и кинетика роста плёнок из адамантана в~тлеющем разряде}

\author{А.\,М.~Штеренберг, В.\,И.~Зынь, А.\,В.~Шацкий,  А.\,А.~Сафонов}
{Ш\,т\,е\,р\,е\,н\,б\,е\,р\,г~А.\,М.,\ З\,ы\,н\,ь\,~В.\,И.,\ Ш\,а\,ц\,к\,и\,й~А.\,В.,\,С\,а\,ф\,о\,н\,о\,в\,~А.\,А.}

\institute{\samgtu.}

\email{E-mails}{ashter53@mail.ru, zyn37@mail.ru, iskandar2003@mail.ru, pomailer@mail.ru}

\otherinf{\fullauthor{Александр Моисеевич Штеренберг}\regalia{д.ф.-м.н., проф.}\position{заведующий
кафедрой} \dept{каф. общей физики и физики нефтегазового производства} 
\fullauthor{Владислав Иванович Зынь} \regalia{д.ф.-м.н., проф.} \position{профессор}  
\dept{каф. общей физики и физики нефтегазового производства}
\fullauthor{Александр Владимирович Шацкий} \regalia{к.ф.-м.н., доц.} \position{ст. преподаватель} \dept{каф. общей физики и физики нефтегазового производства}  
\fullauthor{Александр Александрович Сафонов} \regalia{к.ф.-м.н.} \position{младший научный  сотрудник} \dept{каф. общей физики и физики нефтегазового
производства} }

\keywords{газоразрядная полимеризация в парах адамантана,
структура полимерных плёнок, кинетика роста плёнок из адамантана}

\workabstract{Описываются особенности структуры полимеров,
синтезированных в тлеющем разряде пониженного давления в парах
адамантана на электродах и подложках под <<плавающим>>
потенциалом. Структура формируемых продуктов обуславливается
соотношением вкладов механизмов поверхностной и объёмной
полимеризации в формирование полимера на поверхности.
Определяющими факторами являются вид разряда и его параметры:
разряд постоянного или переменного тока, давление и плотность тока
разряда; наличие протока газа или его отсутствие в связи с подачей
в зону формирования покрытия и выносом из неё частиц дисперсной
фазы; место расположения  подложки в области разряда и вне неё.
Изучена кинетика роста плёнок из адамантана при различных условиях
разряда на электродах из алюминия и меди. Скорости роста полимеров
менялись в пределах 0,5–1,4 нм/с на аноде, 4,2–7,5 нм/с на
катоде для различных режимов разряда и двух материалов
электродов.}

\engtitle{Structure and growth kinetics of films of adamantane in a glow discharge}

\engauthor{A.\,M.~Shterenberg, V.\,I.~Zyn,
A.\,V.~Shatsky, A.\,A.~Safonov}{S\,h\,t\,e\,r\,e\,n\,b\,e\,r\,g~A.\,M.,
Z\,y\,n~V.\,I., S\,h\,a\,t\,s\,k\,y~A.\,V., S\,a\,f\,o\,n\,o\,v~A.\,A.}


\enginstitute{\engsamgtu.}

\otherenginfm{\fullauthor{Alexandr M. Shterenberg} \regalia{Dr. Sci. (Phys. \& Math.)} \position{Head of Dept.}
\dept{Dept. of General Physics \& Oil and Gas Production}
\fullauthor{Vladislav I. Zyn} \regalia{Dr. Sci. (Phys. \& Math.)} \position{Professor} \dept{Dept. of General Physics \&
Oil and Gas Production} \fullauthor{Alexandr V. Shatsky}
\regalia{Ph.\,D. (Phys. \& Math.)} \position{Senior Teacher}
\dept{Dept. of General Physics \& Oil and Gas Production}
\fullauthor{Alexandr~A.~Safonov} \regalia{Ph.\,D. (Phys. \& Math.)} \position{Junior Research Scientist} \dept{Dept. of General
Physics \& Oil and Gas Production}}

\engabstract{The peculiarities of the structure of polymers
synthesized in a low-pressure glow discharge in vapors of
adamantane on the electrodes and substrates at a ``floating''
potential are described. The structure of the product formed is
conditioned by the ratio of deposits of the mechanisms of the
surface and volume  polymerization in the formation of the polymer
on the surface. The determining factors are the type of discharge
and its parameters ---  direct or alternating current discharge,
pressure and current density, presence of the flow of gas or its
absence in connection with the feed in the zone of formation of
coatings and set out from her particles of the dispersed phase;
the location of the substrate in the area of the discharge and
outside of it. The kinetics of growth of films of adamantane at
different discharge conditions on the electrodes of different
materials --- aluminum and copper was studied. The growth rate of
polymers varied between 0.5--1.4~nm/s at the anode, and 4.2--7.5~nm/s at the cathode for different discharge modes and two
electrode materials.}

\engkeywords{discharge polymerization in vapors of adamantane,
structure of polymer films, the growth kinetics of the films of
adamantane}

\date{23/VI/2012}
\revisiondate{15/VIII/2012}


\maketitle



\Section[n]{Введение}
Плазмохимическая или газоразрядная полимеризация различных
соединений исследуется уже несколько десятилетий и стала
классикой, в которой тем не менее имеется много неясностей
[1--3]. Направление с годами не только не теряет актуальности, но,
напротив, приобретает все новые перспективы в связи с ростом
технических возможностей для реализации процессов и развитием
теоретических представлений. Открываются как новые стороны
процесса синтеза плёнок, так и новые возможности полезных
применений. Плазмохимическая полимеризация относится к сфере
высоких технологий  и прекрасно вписывается в развивающиеся на
наших глазах нанотехнологии, позволяя разрабатывать эффективные
процессы управляемого синтеза наноструктур. Появилась информация
об использовании в качестве мономеров в газоразрядной
полимеризации адамантана и его производных [4--6]. Адамантаны с их
уникальной алмазоподобной молекулярной структурой являются
естественными элементами наномира и относятся к одним из
наиболее интересных и перспективных  веществ.

Для разработки физико"=химической модели формирования полимеров в
тлеющем разряде, синтеза определенных структур с заданными
свойствами и установления требующихся при создании технологических
процессов параметров необходимым является знание закономерностей
формирования пленочных структур из адамантана и кинетики роста
полимеров на различных поверхностях в плазмохимическом реакторе.

Структура плёнок и кинетика их роста определяются видом разряда
(переменного или постоянного тока), параметрами разряда
(давлением, плотностью и частотой тока разряда, наличием протока
газа), местом проведения процесса (в катодном пространстве
разряда, в положительном столбе, на границе разряда, вне зоны
разряда), материалом и геометрией подложек (формой и вертикальным,
горизонтальным или другим их расположением в реакторе, условиями
на поверхности синтеза (потенциалом поверхности, степенью её
чистоты) [\citen{ShterZyn:bib:Yasuda}]. Важным является выявление
как общих для всех используемых веществ и условий, так и
специфических, индивидуальных для адамантана закономерностей,
чему и посвящена данная работа.

\smallskip

\Section{Материалы и методы} На основе созданных установок
плазмохимического синтеза, состоящих из реактора для получения
полимеров, электрической и вакуумной частей, системы напуска и
фильтрации, зондирующей и регистрирующей аппаратуры и
подколпачного оборудования по разработанным методикам
[\citen{ShterZyn:bib:ShterPot, ShterZyn:bib:ShterZyn}] с помощью
электронного и силового зондового микроскопов, кварцевого
резонансного датчика, ИК спектрометра и других приборов были
проведены исследования структуры и  кинетики роста полимеров в
условиях синтеза в тлеющем разряде в парах адамантана
$\text{C}_{10} \text{H}_{16}$ постоянного тока и переменного тока
частотой 1~кГц. Диапазон исходных давлений 10–120~Па, плотности
тока разряда 0,1–15~А/м$^2$, разность потенциалов электродов до
2~кВ. Преимущество отдавалось режимам с однократным напуском
исходного вещества, позволяющим проводить процесс с наименьшим
числом неконтролируемых параметров разряда.

Экспериментально  с помощью оптической и электронной микроскопии и
рассеяния света лазера было установлено, что в процессе тлеющего
разряда в парах адамантана достаточно ярко выражены процессы
объёмной полимеризации, приводящие к формированию частиц
дисперсной фазы (порошка). Их наличие во многом определяет синтез
полимеров на любых поверхностях, а именно электродах (частицы
активно встраиваются в плёнки на катодах и редко  "--- на анодах) и
подложках под <<плавающем>> потенциалом.

Зависимости толщин полимера от времени на аноде и катоде на
постоянном токе или на электродах на переменном определялись по
разработанным методикам [\citen{ShterZyn:bib:Safonov}] с помощью
силового зондового микроскопа или кварцевого резонансного датчика.
Для исследования процессов полимеризации за счет активных частиц
из плазмы использовались подложки и тонкие проволочные зонды под
<<плавающим потенциалом>>. Исследовалось влияние материала
электрода на скорости роста плёнок.

Структура полимеров на электродах и подложках изучалась с помощью
электронной и зондовой микроскопии и ИК-спектроскопии.

\newpage

\Section{Структуры формируемых полимеров}

 {\sl 1.1. Структуры, формируемые в тлеющем разряде постоянного
тока.} 

\noindent На~рис.~1 представлены структуры синтезированных на катоде,
аноде и подложке под <<плавающим потенциалом>> полимеров в разряде
постоянного тока в парах адамантана. Увеличение на фотографиях
анода и катода одинаково. При выращивании плёнок в разряде
постоянного тока структура их на аноде и катоде различна.
Поверхность плёнок на аноде обычно ровная, с редкими включениями
макрочастиц, а на катоде структура однотипна со структурой плёнок,
полученных в НЧ-разряде на электродах. Под <<плавающим>>
потенциалом формирование полимера происходит на месте выпадения
частиц дисперсной фазы из разряда на подложку, находящуюся под
разрядным промежутком в районе катода (рис.~1, {\sl в}).

\begin{figure} [b!]
\centering \small 

\includegraphics[height=5.2cm]{./00PIC/fig1a}~\includegraphics[height=5.2cm]{./00PIC/fig1b}

{\sl a \hspace{.5\textwidth} б }

\smallskip

\begin{tabular}{cp{.44\textwidth}}
\raisebox{-5cm}[1pt][2pt]{\hspace{-.88mm}\includegraphics[width=6.925cm]{./00PIC/fig1c}} & {\footnotesize \vspace{1.5cm} Рис. 1. Поверхности плёнок из адамантана на катоде ({\sl а}\/), аноде ({\sl б}\/) и  подложке под <<плавающим>> потенциалом ({\sl в}\/) после 160 секунд горения тлеющего разряда в~парах адамантана; давление 80~Па, \centerline{плотность тока 6~$\text{А/м}^2$}} \\

\raisebox{-1.4cm}[1pt][2pt]{\small \sl в} &  \\

\vspace{4mm}

\end{tabular}



%\caption{Поверхности плёнок из адамантана на катоде ({\sl а}\/), аноде ({\sl б}\/) и  подложке под <<плавающим>> потенциалом ({\sl в}\/) после 160 секунд горения тлеющего разряда в~парах адамантана; давление 80~Па, плотность тока 6~$\text{А/м}^2$}
%
%\vspace{-5mm}
\end{figure}

Строение полимера на подложке под <<плавающим потенциалом>> во
многом определяется местом её расположения, т.\,к. скорость
образования плёнки непосредственно на поверхности примерно на
порядок меньше, чем на электроде, но сильно зависит от количества
встраивающихся макрочастиц. Внешне поверхность может меняться от
однородной без частиц дисперсной фазы до порошка.

\smallskip

 {\sl 1.2. Структуры, формируемые в тлеющем разряде переменного
тока.} 


\begin{figure}[b!]

\vspace{-3mm}

\includegraphics[height=5.3cm]{./00PIC/fig2}~\includegraphics[height=5.3cm]{./00PIC/fig3}  

\setlength{\tabcolsep}{0.4pt}
\begin{tabular}{p{.46\textwidth}cp{.5\textwidth}}


\footnotesize 
Рис. 2. Микрофотография поверхности плёнки из адамантана на электроде в~разряде переменного тока; частота 1~кГц; \centerline{толщина плёнки 0,7~мкм} &~~~~&
\footnotesize 
Рис. 3. Микрофотография порошка на под- \centerline{ложке под <<плавающим>> потенциалом}


\end{tabular}

\vspace{-1mm}

\end{figure}

\setcounter{figure}{3}

\begin{figure} [b!]
\centering \small 

\includegraphics[height=5cm]{./00PIC/fig4a}~\includegraphics[height=5cm]{./00PIC/fig4b}

{\sl a \hspace{.5\textwidth} б }

\smallskip

\caption{Микрофотографии поверхности полимера из адамантана на подложках, расположенных в районе верхнего ({\sl а}\/) и  нижнего ({\sl б}\/) электродов}
%
%\vspace{-5mm}
\end{figure}


\noindent В~электродном разряде переменного тока один из электродов
попеременно является катодом и анодом. Образование на электродах
более однородных структур связано с упорядоченным встраиванием из
объёма в плёнки под действием электрического поля заряженных
глобул, находящихся в прикатодном пространстве, и с высокой
скоростью плёнкообразования за счет полимеризации на активных
центрах поверхности (рис.~2).



Проток газа, как и в случае разряда постоянного тока, влияет на
подачу в зону формирования покрытия и вынос из этой зоны частиц
дисперсной фазы. 
При переменном токе частицы дисперсной фазы имеют
больше возможностей перезарядиться, и они легче увлекаются потоком
газа, что приводит к более интенсивному выпадению порошка в
отдельных местах реактора. На рис.~3 показано формирование
полимера на месте выпадения частиц дисперсной фазы из потока на
подложке, находящейся в месте интенсивного движения макрочастиц.

На рис.~4 показаны полимеры, синтезированные на вертикально
расположенной подложке под <<плавающим>> потенциалом длиной 0,2~м
при горизонтальном расположении электродов разряда соответственно
у верхнего (слева) и нижнего (справа) краев разряда. Увеличение
числа глобул у нижнего края разряда связано с влиянием гравитации
на дисперсные частицы. Также при выходе за пределы разряда у краев
электродов образуются завихрения потоков, приводящие к возможности
осаждения в этих точках частиц дисперсной фазы.

\smallskip

\Section{Сканирующая зондовая микроскопия для  исследования
полимеризованных в газовом разряде покрытий} Применение
сканирующей зондовой микроскопии для исследования полимеризованных
в газовом разряде покрытий на поверхностях различных материалов
позволяет произвести количественную оценку состояния поверхности.
Для оценки изменения состояния поверхности достаточно исследовать
образец до и после обработки и произвести сравнение статистики
распределения элементов поверхности. При сканировании
использовался метод постоянной высоты зонда, при котором зонд
устанавливается на высоте, отвечающей резонансом амплитуде сигнала
для данного покрытия. Далее зонд перемещается над поверхностью,
регистрируя амплитуду тока, протекающего по зонду в зависимости от
координаты. Таким образом, условно минимальная амплитуда сигнала
соотносится с максимальным расстоянием между зондом и
поверхностью.

На рис.~5 представлены результаты зондирования поверхности плёнок
из адамантана, синтезированных в разряде переменного тока частотой
1 кГц на электродах при давлении 80~Па и плотности тока разряда
3~А/м$^2$~({\sl а}) и 7~А/м$^2$~({\sl б}). Значительно большие
перепады высот при больших плотностях тока (в районе 20--30~нм по
сравнению с 5--9~нм) свидетельствует о роли ионной бомбардировки
поверхности.

Изучались и процессы выпадения макрочастиц на подложки из
различных материалов, расположенных в разных местах пространства
реактора в одинаковых условиях. Установлено, что частицы хорошо
оседают на металлические подложки, особенно из алюминия, и плохо
"--- на подложки из углерода.

\smallskip

\Section{Синтез плёнок и порошков вне разряда} Макроструктура
плёнок, выращенных под <<плавающим>> потенциалом, зависит от их
местоположения в реакционном объёме, которое во многом определяет
количество встраивающихся макрочастиц. Применение сканирующей
зондовой микроскопии при исследовании покрытий на поверхностях
различных материалов позволяет произвести количественную оценку
изменения состояния поверхности. Для оценки изменения
шероховатости поверхности достаточно исследовать образец до и
после обработки и произвести сравнение статистики распределения
элементов поверхности. На рис.~5, {\sl в} представлена поверхность
стеклянной подложки под <<плавающим>> потенциалом после нанесения
полимера. Изображения по осям соответственно $1000 \times 1000
\times 40$~нм, толщина полимера порядка 400~нм.




\begin{figure} [t!]
\centering \small 

\includegraphics[height=3.6cm]{./00PIC/fig5a}~\includegraphics[height=3.6cm]{./00PIC/fig5b}

{\sl a \hspace{.5\textwidth} б }

\smallskip

\begin{tabular}{cp{.44\textwidth}}
\raisebox{-4cm}[1pt][2pt]{\hspace{-.88mm}\includegraphics[width=6.925cm]{./00PIC/fig5c}} & {\footnotesize \vspace{1.5cm} Рис. 5. Поверхность пленок из адамантана, полученная сканирующим зондо- \centerline{вым микроскопом}} \\

\raisebox{-1.4cm}[1pt][2pt]{\small \sl в} &  \\

\vspace{4mm}

\end{tabular}



\end{figure}

Как уже указывалось, синтез полимера на подложках под <<плавающим
потенциалом>> во многом определяется их местом размещения,
геометрией, потенциалом поверхности, протоком газа, материалом
подложки. В тлеющем разряде формирование полимеров на различных
поверхностях в плазмохимическом реакторе протекает одновременно с
процессами объёмной полимеризации. Возникающие в объёме разряда дисперсные частицы полимера встраиваются в плёнки и влияют на
структуру и свойства формируемых покрытий. Возможно
преимущественное выпадение на подложку порошка. На рис.~6 показана
крупная частица, имеющая явно выраженную фрактальную структуру и
зафиксированная на месте выпадения порошка на медную подложку,
расположенную на месте выхода частиц дисперсной фазы из объёма
разряда в приэлектродном пространстве.

Рассматриваемые глобулы не являются  чисто механическими
включениями, а образуют специфические для полимеризованных в
разряде веществ внедрения, сросшиеся с окружающим их полимером
посредством химических связей.



Отличие собственного строения глобул от строения полимерных плёнок
объясняется рядом причин:
во"=первых, скорость роста плёнки на электроде на порядок выше скорости роста плёнки на поверхности под <<плавающим>>
потенциалом, а именно им и характеризуется поверхность глобул;
во"=вторых, на рост макрочастиц прямое воздействие оказывают
процессы коагуляции, и скорость роста частиц порошка может
оказаться сравнимой со скоростью роста плёнок и даже превышающей её.


\begin{figure}[t!]



\includegraphics[height=4.8cm]{./00PIC/fig6}~\includegraphics[height=4.8cm]{./00PIC/shter7}  

\setlength{\tabcolsep}{0.4pt}
\begin{tabular}{p{.46\textwidth}cp{.5\textwidth}}


\footnotesize 
Рис. 6. Агрегат дисперсных частиц, выпавший из объёма разряда в парах адамантана на подложку под «плавающим \centerline{потенциалом»} &~~~~&
\footnotesize 
Рис. 7. Кинетика роста полимера в~парах адамантана на алюминиевом и медном электроде в~разряде постоянного тока; давление 100~Па, плотность тока разряда~3 $\text{А/м}^2$; катод: 1 "--- алюминий, 2 "--- медь; анод: 3 "--- алюминий, \centerline{4 "--- медь}


\end{tabular}

\vspace{-8mm}

\end{figure}

Макроструктура плёнок, выращенных под <<плавающим>> потенциалом,
зависит от множества вышеуказанных факторов. Она может меняться от
однородных плёнок без включений макрочастиц до порошков,
полученных в тлеющем разряде с их физико-химической спецификой,
отражающей как процессы коагуляции в объёме, так и химической
сшивки на молекулярном уровне отдельных элементов во всем
диапазоне размеров "--- от  фрагментов исходных молекул до крупных
глобул.

\smallskip



\Section{Кинетика роста плёнок} Изучение кинетики роста плёнок из
адамантана проводилось при различных условиях разряда на
электродах из различных материалов "--- алюминия и меди. На рис.~7
представлена кинетика роста адамантана на алюминиевых и медных
анодах и катодах в разряде постоянного тока при давлении 100~Па и
плотности тока разряда 3~А/м$^2$. Скорости роста полимеров для
разряда в парах  адамантана менялись в пределах 0,5–1,4 нм/с на
аноде, 4,2–7,5 нм/с на катоде для различных режимов разряда.

\smallskip

\Section{Обсуждение} Формирование полимера в тлеющем разряде
происходит как в его объёме, так и на различных поверхностях
подложек, расположенных в плазмохимическом реакторе. При этом
формально можно выделить три возможных варианта: 
\begin{enumerate}
\item 
формирование различных агрегатов дисперсных частиц в объёме
реактора и выпадение их на подложку; 
\item рост плёнки из
активных частиц, молекул мономера с образованием полимерных
островков"=кластеров на поверхности подложки  с дальнейшим
образованием и ростом сплошной плёнки; 
\item образование полимера
в виде композита из формирующейся на поверхности подложки плёнки с
внедрёнными в неё из объёма дисперсными частицами.
\end{enumerate}

Структура композита определяется соотношением скоростей роста
полимера, определяемых механизмами поверхностной и объёмной
полимеризации. Структуры могут меняться от формируемого полимера
практически без участия дисперсных частиц (рис. 8, {\sl а\/}), с
участием плазмазолей (рис. 8, {\sl б}\/) и практически образование
порошка из одних золей (рис. 8, {\sl в}\/). Вид формируемого
композита определяется соотношением скоростей механизмов
поверхностной и объёмной полимеризации. Вклад каждого из процессов
(синтез за счет непосредственного роста на поверхности и за счет
встраивания частиц из объёма) определяется типом разряда, исходным
веществом и многими факторами, связанными с используемой подложкой
(геометрией подложки и реактора, местом  расположения и материалом
подложки, потенциалом её поверхности и другими условиями).




Таким образом, установлено, что структура формируемых продуктов
определяется соотношением вкладов механизмов поверхностной и
объёмной полимеризации в формирование полимера на поверхности. Это
соотношение определяется физико"=химическими параметрами процесса.
Определяющими факторами являются вид разряда и его параметры –
разряд постоянного или переменного тока, давление и плотность тока
разряда; наличие протока газа или его отсутствие в связи с подачей
в зону формирования покрытия и выносом из неё частиц дисперсной
фазы; компоновка реактора, то есть выбор места расположения
подложки в области разряда или вне неё; горизонтальное,
вертикальное или любое другое её размещение.

В результате синтеза образуется 3D-макромолекула, <<сшитая>> в
различных направлениях разными видами соединений. Отдельные
структурные элементы полимера отличаются друг от друга по двум
причинам.  Первая "--- из"=за протекающих в разряде процессов
фрагментации исходных и синтезируемых соединений и постоянной
рекомбинации различных фрагментов различаются как молекулы,
образующие полимер, так и молекулы, их <<сшивающие>>, если вообще
можно делать различия между этими видами молекул в
плазменно"=полимеризированном продукте. Вторая "--- из-за влияния
вышеуказанного соотношения вкладов механизмов поверхностной и
объёмной полимеризации в формирование депозита. Внедрение в плёнки
частиц порошка отражается на макроструктуре полимера.

Структура полимера на электроде зависит от исходного вещества,
времени полимеризации, материала и степени чистоты обработки
поверхности электрода, от условий полимеризации. Уже начиная с~толщины,
большей 0,1~мкм, на поверхности видны глобулярные образования,
которые постепенно укрупняются, а свыше толщины порядка 0,4~мкм вид
структуры плёнок практически не меняется. При большой толщине
структура плёнок мало зависит от материала электродов.

При выращивании плёнок в разряде постоянного тока структура их на
аноде и катоде различна. Поверхность  плёнок на аноде ровная, с
редкими включениями макрочастиц, а на катоде структура более
сложная и содержит признаки радиационных повреждений. Такого же
типа структура характерна для плёнок, полученных в НЧ-разряде на
электродах. В случае НЧ-разряда каждый из электродов в течение
около половины времени полимеризации является катодом, что
приводит к формированию плёнки с сильными признаками участия
макрочастиц и ионной бомбардировки на фоне слабой и фактически
незаметной анодной структуры.

\setcounter{figure}{7}

\begin{figure}[t!]


\small \sl
\centering

\includegraphics[width=\textwidth]{./00PIC/shter8}

а \hspace{.33\textwidth} б \hspace{.33\textwidth} в

\caption{Структура полимера, синтезируемого на подложке при различных скоростях роста полимера, определяемых механизмами поверхностной и объемной полимеризации}

\vspace{-5mm}

\end{figure}

Результаты ИК-спектрометрии показывают, что полосы поглощения для
синтезируемых соединений по отношению ко всем исходным веществам в
результате полимеризации расширяются и несколько смещаются.
Расширение полос говорит о том, что полимеризация в плазме
тлеющего разряда приводит к образованию сшитых структур.

Основные кинетические закономерности, отмеченные по результатам
экспериментов, являются практически общими для всех режимов, 
отличаясь лишь количественно для разных параметров давления и плотности
тока. Во всех режимах после некоторого периода роста наблюдается
тенденция к замедлению и даже насыщению роста массы как на
электродах, так и на изолированных подложках. В то же время для
разряда на постоянном токе скорости роста и масса выросшей плёнки
на катоде, аноде и подложке значительно отличаются. Наибольшие
скорости (на медном катоде до 7,5~нм/с) и массы (толщина
синтезированных плёнок до 2--2,5~мкм за 500~с) обеспечиваются на
поверхности катода, наименьшие "--- на поверхности стеклянной
подложки, где эти показатели на порядок величины меньше. Для анода
получены значения, лишь ненамного большие тех, которые характерны для
подложки. Материал электрода не оказывает значительного влияние на
кинетические процессы, так как основные измерения уже производятся
на пленках значительной толщины (1--2 мкм).

На переменном токе кинетические кривые более резки, подъем более
крутой, и состояние, похожее на насыщение роста, достигается
быстрее на временах порядка 250--350~с. Изменение фрактальной
размерности поверхности плёнок адамантана от времени выращивания
при трех разных плотностях тока, определяемое по методу, предложенному
Мандельбротом [\citen{ShterZyn:bib:Safonov}], демонстрирует
наличие кинетических структурных зависимостей строения
синтезируемых полимеров.

Таким образом, установленные экспериментально закономерности
демонстрируют однозначное влияние параметров разряда на структуру
полимеров, синтезированных в газовом разряде в парах адамантана, и
наличие в их строении кинетических структурных зависимостей.


\thanks{Работа выполнена при поддержке Федеральной целевой
программы <<Научные и научно"=педагогические кадры инновационной
России>> на 2009--2013 годы.}


\begin{thebibliography}{8}

\RBibitem{ShterZyn:bib:TkachukKolot} 
\by Ткачук~Б.\,В., Колотыркин~В.\,М., 
\book Получение тонких полимерных плёнок из газовой фазы 
\publaddr М. 
\publ Химия 
\yr 1977 
\totalpages 321
\misctransl
\by  Tkachuk~B.\,V., Kolotyrkin~V.\,M.
\book Thin polymer film deposition from the gas phase
\publaddr Moscow 
\publ Khimiya
\yr 1977
\totalpages 321 

\Bibitem{ShterZyn:bib:Yasuda} 
\by Yasuda~H.
\book Plasma polymerization
\publaddr New York
\publ Academic Press
\yr 1985
\totalpages 240
\rtransl русск. пер.:~
\by Ясуда~Х., 
\book Полимеризация в плазме 
\publaddr М. 
\publ Мир 
\yr 1988 
\totalpages 376

\RBibitem{ShterZyn:bib:ShterPot}
\book Макрокинетика формирования дисперсной фазы в~газоразрядных системах
\by Штеренберг~А.\,М., Потапов~В.\,К.
\yr 1997
\publ Самара
\publaddr СамГТУ
\totalpages 180
\misctransl
\by Shterenberg~A.\,M., Potapov~V.\,K.
\book Macrokinetics of Dispersed-Phase Formation in Gas Discharge Systems
\publaddr Samara
\publ SamGTU
\yr 1997
\totalpages 180

\RBibitem{ShterZyn:bib:Fedoss}
\paper Полимеризация паров адамантана в электрическом разряде
\by Федосеев~Д.\,В., Лаврентьев~А.\,В., Варшавская~И.\,Г.
\jour Докл. АН СССР
\yr 1986
\vol 291 
\issue 4
\pages 917--919
\misctransl
\paper Polymerization of Adamantane Vapors in Electrical Discharge
\by Fedoseev~D.\,V., Lavrentiev~A.\,V., Varshavskaya~I.\,G.
\jour Dokl. AN SSSR
\yr 1986
\vol 291 
\issue 4
\pages 917--919


\Bibitem{ShterZyn:bib:Carlson}
\preprint Bias enhanced nucleation of diamond films in a chemical vapor deposition process
\by     Dahl~J.\,E., Carlson~R.\,M., Liu~S., Queshi~W.\,R., Bokhari~W.
\preprintinfo US Patent Application Publication; Pub. No.: US 2006/0228479~A1; Pub.  Date: Oct.~12, 2006


\Bibitem{ShterZyn:bib:CarlsonLiu}
\preprint Nucleation of diamond films from higher diamondoids
\by     Dahl~J.\,E., Carlson~R.\,M., Liu~S.
\preprintinfo US Patent Application Publication; Pub. No.: US 2005/0019576~A1; Pub.  Date: Jan.~27, 2005

\RBibitem{ShterZyn:bib:ShterZyn}
\paper Полимеризация паров адамантана в тлеющем разряде
\by Андреева~А.\,В., Зынь~В.\,И., Сафонов~А.\,А., Шацкий~А.\,В., Штеренберг~А.\,М.
\jour  Извест. Самарск. научн. центра Российской академии наук
\yr 2011
\vol 13
\issue 4
\pages 84--90
%http://elibrary.ru/item.asp?id=17304552
\misctransl
\paper  Polymerization of Adamantane Vapors in Glow  Discharge
\by Andreeva~A.\,V., Zyn~V.\,I., Safonov~A.\,A.,  Shatsky~A.\,V., Shterenberg~A.\,M.
\jour  Izvest. Samarsk. Nauchn. Centra Rossiyskoy  Akademii Nauk 
\yr 2011
\vol 13
\issue 4
\pages 84--90



\RBibitem{ShterZyn:bib:Safonov} 
\by Сафонов~А.\,А. 
\thesis Кинетика формирования и свойства нано- и микроструктур полимеров, синтезируемых в разряде пониженного давления в парах адамантана
\thesisinfo Дисс. канд. \dots\ физ.-мат. наук: 01.04.17~-- Химическая физика, горение и взрыв, физика экстремальных состояний вещества 
\publaddr Самара 
\yr 2012 
\totalpages 152
\misctransl
\by Safonov~A.\,A.
\thesis The kinetics of formation and properties of nano- and microstructures of the polymers synthesized in the discharge of low pressure in the vapor adamantane
\thesisinfo Ph.D. Thesis  (Phys. \& Math.)
\publaddr Samara
\yr 2012 
\totalpages 152

\end{thebibliography}


\noindent
\hskip6.5cm \thedate \par\noindent
\hskip6.5cm\therevdate



%\newpage
%
%\chead[\fancyplain{}{\scriptsize \sl S\,h\,t\,e\,r\,e\,n\,b\,r\,g~A.\,M., Z\,y\,n\,~V.\,I.,
%S\,h\,a\,t\,s\,k\,y~A.\,V., S\,a\,f\,o\,n\,o\,v~A.\,A.}]{\fancyplain{}{\scriptsize \sl  Structure and growth kinetics of films from adamantane in a glow discharge \dots}}

\makeengtitlem

\renewcommand{\baselinestretch}{0.9}

\newpage

%
% \tableofengcontents
%
%\end{document}
